\documentclass[11pt]{article}

\usepackage[margin=2.6cm]{geometry}
\usepackage{amsmath,amssymb}
\usepackage{booktabs}
\usepackage{array}
\usepackage{xcolor}
\usepackage{hyperref}
\usepackage{enumitem}
\usepackage[T1]{fontenc}
\usepackage{lmodern}
\usepackage{parskip}

\hypersetup{
    colorlinks=true,
    linkcolor=blue!50!black,
    citecolor=blue!50!black,
    urlcolor=blue!50!black
}

\newcommand{\code}[1]{\texttt{#1}}

\title{\vspace{-1.5cm}Beyond Nested Sampling: A Survey of Samplers for\\ Multimessenger Bayesian Inference\\[4pt]
\large Research notes for the NMMA / retreat discussion}
\author{Prepared with Claude Code (research assistant), for Thibeau Wouters}
\date{September 2026}

\begin{document}
\maketitle

\begin{abstract}
\noindent
NMMA performs joint Bayesian inference over gravitational-wave (GW), electromagnetic
(kilonova/GRB) and nuclear equation-of-state (EOS) data by wrapping \code{bilby} around
nested samplers (\code{pymultinest} by default, plus \code{dynesty} and \code{ultranest}).
This note surveys (i) which other codes in the multimessenger/nuclear-astrophysics space
solve structurally similar problems, whether they are open source, and which samplers they
use; (ii) the \emph{full} menu of samplers already reachable through \code{bilby} without
any new integration work, most of which NMMA does not currently exercise; (iii) structural
alternatives that change \emph{what} gets sampled rather than only how; and (iv) what
neighbouring and more distant Bayesian-computation communities facing high-dimensional,
expensive-likelihood inference have already adopted. Every performance number is checked
directly against each paper's own \LaTeX{} source rather than quoted from an abstract.
Headline finding: essentially the entire multimessenger/nuclear-EOS ecosystem (NMMA, bajes,
redback, NEoST, the nonparametric-EOS lineage) still relies on the same CPU-bound
nested-sampling engines, while the pure-GW subfield next door has already moved to
GPU-native, gradient-based, amortized, or surrogate-based methods with verified $5\times$ to
$\sim\!400\times$ wall-clock reductions (and, for Gaussian-process surrogates on expensive
likelihoods, up to two orders of magnitude \emph{fewer likelihood evaluations}) on problems
of comparable or higher dimensionality --- including two pieces of work the group itself has
already co-authored. This is a concrete, low-risk opportunity, not a hypothetical one.
\end{abstract}

\section{The starting point: how NMMA samples today}

A direct inspection of the NMMA source
(\url{https://github.com/nuclear-multimessenger-astronomy/nmma}) confirms the following,
rather than assuming it from the papers alone:

\begin{itemize}
\item NMMA's inference stack is built on \code{bilby}~\cite{Ashton:2018jfp}, following the
original and updated NMMA methodology papers~\cite{Coughlin:2018fis,Pang:2022rzc}, with an
extended framework paper released this year~\cite{Rose:2026cnm}. All of its supported
samplers are CPU-bound nested samplers: \code{pymultinest} is the CLI default (MultiNest,
\cite{Feroz:2008xx}), \code{dynesty}~\cite{Speagle:2019ivv} is a core dependency, and
\code{ultranest}~\cite{Buchner:2021cql} ships as an optional extra
(\code{nmma/core/parsing.py}, \code{pyproject.toml}).
\item EOS inference solves the Tolman--Oppenheimer--Volkoff equations directly inside the
likelihood: \code{nmma/eos/tov.py} integrates the stellar-structure ODEs with
\code{scipy.integrate.solve\_ivp} for \emph{every} sampled EOS, at every mass on the mass
grid, on every likelihood call --- no caching, surrogate, or amortization across samples.
\item Parallelism is handled purely at the level of MPI across independent likelihood
evaluations (\code{nmma/core/mpi\_setup.py}): more CPU workers, not GPU or
gradient-based acceleration.
\end{itemize}

This is a reasonable default: nested sampling is robust to multimodality, needs no
gradients, and directly returns the Bayesian evidence NMMA needs for model selection between
EOS families or kilonova/GRB models. The cost is that nested sampling's live-point evolution
requires $\mathcal{O}(10^5\text{--}10^7)$ likelihood evaluations per run, growing with
dimensionality and with the correlation length of the target. When the likelihood is itself
slow --- GW waveform generation, an on-the-fly TOV solve, a kilonova/GRB light-curve
evaluation --- the sampler's evaluation count directly sets the wall-clock budget, and joint
multimessenger analyses stack several such slow likelihoods (and their parameters) at once.
This compounds exactly where NMMA's science ambitions (more sources, more parameters,
next-generation-detector sensitivities) are pushing hardest.

\section{The full sampler menu: what's already one flag away}
\label{sec:menu}

NMMA declares only three samplers as dependencies, but it inherits \emph{all} of
\code{bilby}'s sampler backends for free, since NMMA's likelihoods are ordinary
\code{bilby.Likelihood} objects. Querying the installed \code{bilby} directly
(\code{bilby.core.sampler.get\_implemented\_samplers()}) rather than reading documentation
gives the authoritative list in Table~\ref{tab:menu}: 16 samplers, of which NMMA's own
documentation surfaces only 3.

\begin{table}[htbp]
\centering
\small
\begin{tabular}{@{}p{2.6cm}p{2.3cm}p{9.0cm}@{}}
\toprule
\textbf{Sampler} & \textbf{Family} & \textbf{Notes} \\
\midrule
pymultinest & Nested & NMMA's current default~\cite{Feroz:2008xx}. \\
dynesty & Nested (dynamic) & NMMA core dependency~\cite{Speagle:2019ivv}. \\
ultranest & Nested (reactive) & NMMA optional extra~\cite{Buchner:2021cql}. \\
nestle & Nested & Lightweight dynesty precursor; no dedicated paper. \\
cpnest & Nested & Used by earlier \code{bajes} versions~\cite{Breschi:2021wzr}. \\
dnest4 & Nested (diffusive) & Handles strongly multimodal posteriors. \\
pypolychord & Nested (slice) & Scales better than rejection-based nested sampling in high dimensions~\cite{Handley:2015vkr}. \\
nessai & Nested + normalizing flow & Flow-informed proposals~\cite{Williams:2021qyt,Williams:2023ppp}. \\
emcee & MCMC (ensemble, affine-invariant) & Used by MOSFiT~\cite{Guillochon:2017bmg}. \\
zeus & MCMC (ensemble slice) & Gradient-free, tuning-free, efficient on curved/correlated posteriors~\cite{Karamanis:2021tsx}. \\
kombine & MCMC (kernel-density ensemble) & No dedicated paper found. \\
ptemcee & MCMC (parallel-tempered ensemble) & Better multimodal mixing than plain emcee. \\
ptmcmcsampler & MCMC (parallel-tempered) & LALInference-lineage PTMCMC. \\
bilby\_mcmc & MCMC (native, block-proposal) & Native to bilby~\cite{Ashton:2021anp}; see below. \\
pymc & \textbf{Gradient-based (HMC/NUTS)} & Only entry point to gradient-based sampling already in the menu. \\
\bottomrule
\end{tabular}
\caption{All samplers exposed by \code{bilby.core.sampler.get\_implemented\_samplers()} in
the installed environment (excluding the internal \code{fake\_sampler} test stub). Verified
directly from the installed package, not from documentation.}
\label{tab:menu}
\end{table}

Two things stand out. First, \textbf{\code{pymc} is already in the menu}: it is bilby's only
built-in route to gradient-based (HMC/NUTS) sampling. It cannot help NMMA today because
NMMA's TOV solve and waveform generation are not differentiable end-to-end, but the
integration point already exists --- differentiability, not sampler plumbing, is the
remaining blocker (see Section~\ref{sec:crosscult2}). Second, \textbf{\code{zeus}} and
\textbf{\code{bilby\_mcmc}} are installable/available today with no code change beyond the
\code{--sampler} string, yet neither appears in NMMA's examples or dependency list.
\code{bilby\_mcmc}'s own validation paper is worth reading before adopting it, however: it
reports that dynesty remains an order of magnitude more sample-efficient per core than
\code{bilby\_mcmc}, and the latter's advantage is specifically that it is embarrassingly
parallelizable across many independent short chains in a high-throughput-computing
environment, not that any single chain is faster~\cite{Ashton:2021anp}. It is a genuine
option for cluster-heavy workflows, not a free lunch.

\section{Who else samples multimessenger sources, and how?}

Searching INSPIRE-HEP for codes and papers performing Bayesian parameter estimation on GW,
EM-transient, or EOS data in a way structurally comparable to NMMA turns up a small,
well-defined set of tools, summarized in Table~\ref{tab:codes}.

\begin{table}[htbp]
\centering
\small
\begin{tabular}{@{}p{2.3cm}p{2.0cm}p{1.1cm}p{3.9cm}p{3.9cm}@{}}
\toprule
\textbf{Code} & \textbf{Reference} & \textbf{Open?} & \textbf{Sampler(s)} & \textbf{Scope} \\
\midrule
NMMA & \cite{Coughlin:2018fis,Pang:2022rzc,Rose:2026cnm} & Y & pymultinest, dynesty, ultranest (via bilby) & Joint GW+EM+EOS \\
bajes & \cite{Breschi:2021wzr,Breschi:2024qlc} & Y & dynesty (previously cpnest) & Joint GW+kilonova \\
NEoST & \cite{Raaijmakers:2025hbz} & Y & Nested sampling (MultiNest/dynesty-family) & EOS from NICER+GW \\
Nonparametric EOS (GP/spectral) & \cite{Landry:2020vaw,Legred:2021hdx,Finch:2025bao,Ng:2025wdj} & Mixed & Nested/importance/rejection sampling over GP-drawn EOS priors & EOS from GW(+EM) \\
redback & \cite{Sarin:2023khf} & Y & dynesty (default), pymultinest, nestle (via bilby, 16 samplers total) & EM transients only (kilonova, GRB, SN) \\
MOSFiT & \cite{Guillochon:2017bmg} & Y & emcee (ensemble MCMC) & EM transients only \\
\midrule
jim / flowMC & \cite{Wong:2023lgb,Wong:2022xvh} & Y & Normalizing-flow-enhanced MCMC (JAX/GPU) & GW only \\
GPU nested sampling & \cite{Yallup:2026nqc} & Y & Nested sampling, GPU-vectorized Slice-within-Gibbs & GW only \\
dingo & \cite{Dax:2021tsq,Dax:2022pxd} & Y & Amortized neural posterior estimation & GW only \\
nessai & \cite{Williams:2021qyt,Williams:2023ppp} & Y & Nested sampling with normalizing-flow proposals & GW only \\
\bottomrule
\end{tabular}
\caption{Codes performing Bayesian inference on multimessenger-adjacent sources. The top
block is directly comparable to NMMA's problem class (multi-channel or EOS inference); the
bottom block is GW-only but included because it is the same likelihood family (compact-binary
waveform models) and shows what happens once nested sampling is accelerated or replaced.}
\label{tab:codes}
\end{table}

Two observations stand out. First, \textbf{bajes}~\cite{Breschi:2021wzr,Breschi:2024qlc} is
the closest thing to a genuine NMMA competitor: it performs joint, coherent GW+kilonova
inference, is open source, and --- checked directly in the source of its 2024 joint-inference
paper --- uses the \code{dynesty} nested sampler (having used \code{cpnest} in earlier work).
No open-source code was found that performs joint GW+EM+\emph{EOS} inference together the way
NMMA does; NEoST and the Landry/Essick/Legred/Chatziioannou nonparametric-EOS lineage handle
the EOS side alone. Second, and more importantly, \textbf{every one of these
directly-comparable codes is built on nested sampling}: dynesty, MultiNest, cpnest, or
nestle. The multimessenger and nuclear-EOS subfield has not yet adopted any of the
alternatives described in Section~\ref{sec:crosscult}, even though the pure-GW subfield next
door already has.

\section{Where the neighbouring GW community has already moved, with numbers}
\label{sec:crosscult}

The most relevant lesson does not require leaving gravitational-wave astronomy at all: GW
parameter estimation has the same likelihood structure as NMMA's GW branch (an expensive,
smooth, differentiable-in-principle waveform likelihood over $\sim$10--15 parameters), and it
has already been substantially accelerated away from plain CPU-bound nested sampling. Every
number below was checked directly against the paper's own \LaTeX{} source (fetched via
\code{arxiv.org/e-print} and grepped), not taken from an abstract.

\begin{itemize}[leftmargin=1.4em]
\item \textbf{jim / flowMC}~\cite{Wong:2023lgb,Wong:2022xvh}: a JAX-based, GPU-native,
normalizing-flow-enhanced MCMC sampler. Verified in the jim paper's source: for GW150914 and
GW170817, jim produces $\sim$2500 and $\sim$3500 effective samples respectively, with a total
wall time on a single NVIDIA A100 GPU of \textbf{$\sim$3 minutes}, of which only $\sim$40\,s
is the actual sampling (the rest is one-off JIT-compilation overhead). The same paper reports
the equivalent \code{parallel-bilby}~\cite{Talbot:2019okv}+\code{dynesty} run, executed on
\textbf{400 Intel Skylake CPU cores}, taking $\sim$2\,h for GW150914 and $\sim$1\,day for
GW170817. That is a $\gtrsim\!40\times$ to $\gtrsim\!400\times$ wall-clock reduction using one
GPU instead of hundreds of CPU cores, on the same data and priors.
\item \textbf{GPU-native nested sampling}~\cite{Yallup:2026nqc} (co-authored within the
group): rather than replacing nested sampling, this keeps the algorithm and its evidence
output exactly as-is and instead re-implements its inner Slice-within-Gibbs kernel to run
vectorized on GPU, batching waveform generation across chains. Verified directly in the
paper's abstract and body: for \emph{full-duration, uncompressed} long binary-neutron-star
signals across a three-detector network --- the exact signal class NMMA's GW branch targets
--- it delivers well-calibrated posteriors in a \textbf{median of twelve minutes on a single
GPU} (five minutes sharded across four GPUs), with no approximation to the likelihood. Adding
heterodyning/relative-binning compression~\cite{Zackay:2018qdy,Dai:2018dca} brings the median
down to \textbf{89 seconds}, and a precessing-spin, tidal re-analysis of GW170817 itself
completes in \textbf{$\sim$2 minutes}. Because this changes nothing about how NMMA's
scientists would interpret the output --- still nested sampling, still a Bayesian evidence,
still the full physical waveform --- it is arguably the lowest conceptual-risk transplant on
this list.
\item \textbf{Evidence estimation with jim + learned harmonic mean}~\cite{Polanska:2024zpn}
(also co-authored within the group): benchmarks GPU-based normalizing-flow evidence
estimation (jim + \code{harmonic}, 1 GPU) against \code{parallel-bilby}+dynesty (16 CPU
cores) on a 4-dimensional and an 11-dimensional GW inference problem. Verified in the paper's
own source: consistent evidence estimates with speedups of \textbf{$5.4\times$} (4D) and
\textbf{$14.8\times$} (11D) --- the speedup \emph{grows} with dimensionality, the opposite of
nested sampling's usual scaling, and exactly the regime NMMA's joint high-dimensional
parametrizations live in.
\item \textbf{dingo}~\cite{Dax:2021tsq,Dax:2022pxd}: amortized neural posterior estimation
--- pay a one-off training cost (verified in the PRL paper's source: \emph{``training takes
roughly 10 days''} on a GPU), then obtain a full posterior for a \emph{new} event in seconds.
The paper's own text states inference time is reduced \emph{``from $\mathcal{O}(\text{day})$
to 20 seconds per event,''} cross-validated against \code{bilby}+\code{dynesty} and
LALInference nested-sampling posteriors via Jensen--Shannon divergence.
\item \textbf{nessai}~\cite{Williams:2021qyt,Williams:2023ppp} takes a hybrid route: keep
nested sampling but replace the inefficient rejection-based live-point proposal with a
normalizing-flow-informed one, staying entirely inside the framework NMMA already depends on.
\end{itemize}

The common thread: a smooth, repeatedly-evaluated, GPU-batchable, or learnable likelihood
surface is exactly what plain nested sampling never exploits, and what GPU-vectorized nested
sampling, flow-enhanced MCMC, amortized neural posterior estimation, and flow-informed
proposals all do, in four increasingly aggressive ways.

\section{Structural alternatives: change what gets sampled, not just how}
\label{sec:structural}

Everything above still evaluates the same likelihood function the same number of
conceptual times, just faster per evaluation (GPU) or with fewer wasted evaluations (flows).
A different, complementary line of work instead restructures the inference problem itself so
that far \emph{fewer} raw likelihood evaluations are needed in the first place --- directly
attacking the cost of NMMA's on-the-fly TOV solve rather than only the cost of the sampler
wrapped around it.

\begin{itemize}[leftmargin=1.4em]
\item \textbf{GPry}~\cite{Gammal:2022eob}: builds an active-learning Gaussian-process
surrogate of the log-posterior itself, using a support-vector-machine classifier to reject
non-finite regions, and queries the true (expensive) likelihood only at a small, sequentially
chosen set of points. Verified directly in the paper's abstract and body: it \emph{``reduces
the number of required posterior evaluations by two orders of magnitude compared to
traditional Monte Carlo inference,''} demonstrated concretely on a 6-dimensional Planck CMB
likelihood converging in $\sim$500 evaluations of the underlying Boltzmann code (versus
$\sim\!10^4$--$10^5$ for MCMC/nested sampling), and on a 4-dimensional BBN+BAO likelihood
converging in 80--124 evaluations. Crucially, \textbf{it needs no GPU and no pre-training}
--- it is a pip-installable, drop-in replacement for the sampler, and the paper explicitly
targets exactly NMMA's regime: \emph{``enabling inference from extremely slow likelihoods
($\gtrsim$ minutes per evaluation), which otherwise would be impossible to sample.''} This is
arguably the most direct match to the TOV-solve cost problem of anything surveyed in this
document.
\item \textbf{RIFT}~\cite{Lange:2018pyp}: rather than treating all parameters uniformly,
RIFT splits them into cheap ``extrinsic'' parameters (analytically/semi-analytically
marginalized for a given point) and expensive ``intrinsic'' parameters, then builds an
iteratively-refined Gaussian-process interpolant of the marginalized likelihood
\emph{over the intrinsic parameters only}, and samples from that interpolant instead of the
raw likelihood. Verified in the paper's own method section: this is a deliberate
generalization of Pankow et al.'s grid-interpolation approach to an unstructured,
actively-refined grid. RIFT is open source and used in production LIGO/Virgo/KAGRA parameter
estimation. \textbf{The structural analogy to NMMA is direct}: NMMA's EOS parameters are
exactly an ``expensive intrinsic'' block (each draw costs a fresh TOV solve), while other
parameters are comparatively cheap --- a RIFT-style iteratively-refined GP interpolant over
EOS-parameter space, built once and re-queried, could replace repeated raw TOV solves during
sampling.
\item \textbf{dot-PE / cogwheel}~\cite{Mushkin:2025yks,Roulet:2022kot,Roulet:2024hwz}: the
most radical option, and verified as accurately described: for compact-binary GW inference,
this line of work (relative-binning plus pre-selected intrinsic/extrinsic/phase parameter
grids) evaluates the likelihood for \emph{all} grid combinations simultaneously via dense
matrix multiplication, dispensing with a stochastic sampler entirely. The 2025 paper reports
full parameter estimation and evidence computation, cross-validated against
\code{cogwheel}+\code{nautilus}~\cite{Lange:2023ydq} nested sampling, in \textbf{a few
minutes on a single CPU core}. This is the least directly transferable to NMMA's EOS/EM
branches (it relies on the specific bilinear structure of compact-binary matched-filtering
likelihoods), but it is a useful reminder that ``sample less'' is not the only lever ---
``don't sample at all, just evaluate everywhere cheaply'' is sometimes possible when the
likelihood has exploitable algebraic structure.
\end{itemize}

\section{Learning from communities further away}
\label{sec:crosscult2}

The multimessenger/nuclear-astro community is not alone in having started from
MCMC/nested-sampling defaults built for modest dimensionality and cheap-ish likelihoods.
Other communities solved essentially the same scaling problem earlier, and each embodies a
different strategy that maps onto a different piece of NMMA's likelihood stack.

\begin{description}[leftmargin=1.2em,style=nextline]
\item[Cosmology --- surrogate the slow physics, not (only) the sampler.]
CMB/large-scale-structure inference (\code{Cobaya}~\cite{Torrado:2020dgo},
\code{CosmoSIS}~\cite{Zuntz:2014csq}) faces an expensive deterministic solve (a Boltzmann
code) inside every likelihood evaluation --- the direct analogue of NMMA's on-the-fly TOV
solve for the EOS, or waveform generation for GW. The community's answer was largely
\emph{not} a fancier sampler but a fast, differentiable emulator substituted for the slow
solver, after which ordinary gradient-based samplers become affordable. For NMMA, the direct
analogue is investing in JAX-differentiable or neural-emulated TOV/EOS and waveform
evaluations, rather than only chasing sampler-level speedups.
\item[Probabilistic programming --- go gradient-based.] Outside astrophysics,
high-dimensional Bayesian inference is dominated by Hamiltonian Monte Carlo / NUTS via
\code{Stan}, \code{PyMC}, or \code{NumPyro}~\cite{Phan:2019elc}. These scale far better with
dimension than nested sampling because proposal efficiency does not degrade combinatorially
with parameter count, but they require a differentiable log-likelihood. This is the real
prerequisite shared by jim and the emulator route above: a JAX (or otherwise
autodiff-compatible) rewrite of the likelihood is the single change that unlocks HMC/NUTS,
flow-based MCMC, and gradient-informed nested sampling simultaneously.
\item[Simulation-based inference --- amortize across events.] The SBI paradigm that dingo
imports into GW astronomy (train once on simulated parameter$\to$data pairs, infer new
events near-instantly), reviewed broadly in~\cite{Cranmer:2019eaq} and used heavily for LHC
inference via tools such as \code{MadMiner}~\cite{Brehmer:2019xox}, is a natural fit for
NMMA's EM branch: kilonova and GRB afterglow simulators are not free to evaluate, but they
\emph{can} be called offline ahead of time. The repeated per-event nested-sampling cost NMMA
currently pays for every new kilonova/GRB fit is exactly the cost SBI amortizes away.
\end{description}

\section{Recommendations}

\begin{enumerate}[leftmargin=1.4em]
\item \textbf{NMMA is not behind the curve within its own subfield.} No comparable open
multimessenger or nuclear-EOS code (bajes, redback, MOSFiT, NEoST, the nonparametric-EOS
lineage) has moved past nested sampling either. The opportunity is to be first, not to catch
up.
\item \textbf{Cheapest possible experiment, zero new dependencies:} NMMA already inherits 16
samplers through \code{bilby} (Table~\ref{tab:menu}), of which it exercises 3. Before any
larger engineering effort, simply benchmarking \code{--sampler=zeus} or
\code{--sampler=pypolychord} on an existing NMMA analysis costs a config-file change and an
afternoon, and would establish a baseline for how much is available without touching the
likelihood at all.
\item \textbf{Most direct near-term transplant, already in-house:} evaluate the GPU-native
nested-sampling kernel of~\cite{Yallup:2026nqc} on NMMA's GW likelihood. Unlike every other
item on this list, it changes nothing about the inference paradigm NMMA's science depends on
--- still nested sampling, still an evidence estimate, still the full physical waveform ---
only where the likelihood is evaluated. It was demonstrated on exactly NMMA's target signal
class (long-duration, three-detector BNS) and already has an author in common with the group,
making it plausibly the lowest-effort, highest-confidence win on this list.
\item \textbf{Lowest-risk incremental step:} add \code{nessai}~\cite{Williams:2021qyt} as a
sampler backend alongside pymultinest/dynesty/ultranest. It is a change to the sampler
string, not the likelihood code, and stays inside the nested-sampling/evidence framework NMMA
already depends on.
\item \textbf{Best match to the TOV-solve cost specifically, and probably underrated:} trial
\code{GPry}~\cite{Gammal:2022eob} directly on NMMA's EOS-inference likelihood. It needs no
GPU, no pre-training, and no likelihood rewrite (unlike the JAX-based options below); it is a
pip-installable wrapper. Given its verified two-orders-of-magnitude reduction in required
likelihood evaluations on similarly-dimensioned problems, and that it is explicitly designed
for likelihoods costing seconds-to-minutes per call (exactly NMMA's TOV-solve regime), this
may be the single highest-return-for-effort item on this entire list, and is nearly free to
test.
\item \textbf{Highest-leverage step, also building on in-house work:} extend the jim/flowMC +
learned-harmonic-mean pipeline, already validated for pure GW parameter estimation and
evidence computation by the group itself~\cite{Polanska:2024zpn,Wong:2023lgb}, to NMMA's
joint GW+EM+EOS likelihood. The verified $5.4\times$--$14.8\times$ (evidence, single-GW,
growing with dimension) and $\sim$40$\times$--400$\times$ (full PE, GW150914/GW170817)
speedups, on problems of comparable or lower dimensionality than NMMA's joint parametrization,
suggest large, quantifiable gains rather than aspirational ones.
\item \textbf{Medium-term, highest structural payoff:} pursue differentiable (JAX-based) or
neural-emulated versions of NMMA's two slow deterministic solves --- the GW waveform and the
TOV/EOS integration --- following the cosmology community's emulator strategy. \code{pymc} is
already reachable through \code{bilby} (Table~\ref{tab:menu}), so a differentiable likelihood
would make gradient-based samplers (HMC/NUTS) and amortized SBI available across \emph{all}
of NMMA's source types at once, rather than as piecemeal per-messenger fixes. Alternatively,
a RIFT-style~\cite{Lange:2018pyp} iteratively-refined Gaussian-process interpolant built once
over EOS-parameter space is a lower-effort structural change that targets the same TOV-solve
bottleneck without requiring a differentiable TOV solver at all.
\item \textbf{Longer-term, for the EM branch specifically:} evaluate an amortized/SBI
approach (dingo-style neural posterior estimation) for kilonova and GRB fitting, where the
same light-curve model is currently re-fit from scratch, at full nested-sampling cost, for
every new event.
\end{enumerate}

\section*{Methodology note}
All literature identification and citation metadata in this document were retrieved
programmatically from the INSPIRE-HEP Literature API
(\url{https://inspirehep.net/api/literature}) and are recorded as full BibTeX entries in
\code{references.bib}. Every performance and sampler claim attributed to jim, the GPU-native
nested-sampling kernel, the Polanska et al.\ evidence pipeline, dingo, redback/bajes, GPry,
RIFT, and dot-PE was verified by fetching the paper's own \LaTeX{} source from arXiv
(\code{arxiv.org/e-print/<id>}) and grepping the relevant passages directly, rather than
relying on abstracts. Table~\ref{tab:menu}'s sampler list was obtained by directly querying
the installed \code{bilby} package
(\code{bilby.core.sampler.get\_implemented\_samplers()}), not from documentation. This
process also caught and corrected two errors surfaced during drafting: redback's third
supported sampler is \code{nestle}, not \code{nessai}, and the parallel-bilby citation key is
\code{Talbot:2019okv}, not \code{Smith:2019ucc}. It also surfaced a nuance an uncritical read
would have missed: \code{bilby\_mcmc}'s own validation paper reports it as an order of
magnitude \emph{less} sample-efficient per core than dynesty, not more --- its value is
parallelizability, not raw speed. NMMA's own sampler configuration was verified directly from
its source code (\code{nmma/core/parsing.py}, \code{pyproject.toml}, \code{nmma/eos/tov.py})
rather than assumed from its publications.

\bibliographystyle{unsrt}
\bibliography{references}

\end{document}
